\section{Limitations and Adversarial Cases}
Despite the performance gains observed, the ARS framework is subject to specific architectural limitations and adversarial scenarios where its efficiency may degrade.

\subsection{Bucket Collapse and Pathological Skew}
The Aero architecture relies on a 1024-entry lookup table to approximate the data distribution. However, when the input entropy is extremely low or concentrated in a range smaller than the machine's precision (e.g., millions of elements differing only in the 128th bit), "Bucket Collapse" occurs. In this state, the mapping function $\mathcal{B}(x)$ maps the majority of elements to a single bin, effectively nullifying the spatial partition and forcing the algorithm to revert to its $O(n \log n)$ comparison-based sub-sorter for the entire dataset.

\subsection{The Recursive Refinement Penalty}
While ARS Exp A provides a path to $O(n)$ for arbitrary precision through recursion, each level of recursive mapping incurs a fresh "Sampling Penalty." Constructing the reciprocal multiplier and histogram for every nested bin introduces significant constant-factor overhead. For datasets with high prefix-collisions (e.g., long identical strings), the cost of multiple mapping passes can exceed the cost of a highly optimized $O(n \log n)$ implementation like IPS4o.

\subsection{Memory Overhead and Cache Pressure}
The $O(n)$ space complexity of ARS is its primary trade-off. By requiring an auxiliary scratch buffer and a spatial key cache, ARS consumes approximately $2.5 \times$ more memory than in-place sorters. At massive scales, this increases TLB pressure and can lead to Page Faults if the dataset size approaches the system's physical memory limit, a constraint that comparison-based sorters handle more gracefully.

\subsection{Instability Under Non-Monotonic Projections}
The integrity of the sort depends entirely on the monotonicity of the $\mathcal{K}$ function. If $\mathcal{K}$ is poorly defined or fails to capture the significant bits of a custom data type, the resulting spatial grid will be misaligned, leading to high classification variance. Furthermore, the 8-byte prefix mapping used for strings is blind to characters beyond the first eight bytes during the partitioning phase, making the algorithm sensitive to datasets consisting of long, shared prefixes.

\subsection{Streaming P99 Latency vs. Throughput Trade-off}
In the streaming architecture (Exp D), there is a strict, unavoidable trade-off between peak ingestion throughput and P99 latency. Increasing the micro-batch threshold significantly improves overall throughput but can cause periodic latency spikes during background consolidations. For systems with hard microsecond-scale real-time deadlines, the consolidation phase may induce unacceptable jitter, particularly when processing heavily skewed (Zipfian) data where merging large "heavy-hitter" runs saturates memory bandwidth.
